% ============================================================
\section{Design Discussion}

\subsection{Resources analysis}
PLACEHOLDER

compare the number of DSPs, registers etc.

\subsection{Impact of Hardware Multiplier}

Enabling the Nios~II embedded multiplier reduces latency because the computation of $F(x)$ involves several multiplications per iteration, particularly in the evaluation of $x^3$. 
Although floating-point operations are still emulated in software, these routines internally rely on integer-level mantissa multiplication. 
Providing dedicated hardware for integer multiplication therefore accelerates the most computationally intensive part of the floating-point multiply operation.

The trigonometric term remains software-based. 
Functions such as \texttt{cosf} are implemented using polynomial approximations, which require multiple floating-point operations internally. 
As a result, while hardware multiplier support delivers a clear speedup, the $\cos(\cdot)$ evaluation continues to be a major contributor to total runtime.


\subsection{Floating-Point Hardware Units}

Mapping floating-point addition, subtraction and multiplication to dedicated hardware units further reduces execution time by eliminating the need for full software emulation of IEEE-754 operations. 
In this configuration, the arithmetic core (including mantissa multiplication, exponent alignment, normalisation and rounding) is implemented directly in hardware, significantly shortening the latency of each floating-point operation compared to the software baseline.

However, the Nios~II custom instruction interface is blocking in nature: the processor stalls until the custom instruction asserts completion. 
As a result, even if the floating-point unit itself is internally pipelined, the CPU cannot issue new dependent operations while waiting for the current one to finish. 
This limits effective pipeline utilisation and prevents full throughput exploitation, especially within a sequential loop such as the evaluation of $F(x)$. 
Consequently, while hardware floating-point units reduce per-operation latency, the achievable system-level speedup is constrained by the processor--accelerator interaction model.
